O Ensino a Distância (EAD) pode ser definido como um "sistema tecnológico de comunicação bidirecional, que pode ser massivo, baseado em uma ação sistemática e conjunta de recursos didáticos e o apoio de tutoria, que, separados fisicamente dos estudantes, propiciam a esses uma aprendizagem independente" (Arteiro, 2001, \textit{apud} \cite{rochaMetodologiasAprendizagensNo2022}). Essa modalidade vem crescendo significativamente, sobretudo no ensino superior, com destaque para o fato de que 44\% das matrículas brasileiras em 2020 já eram realizadas nessa forma de ensino \cite{rochaMetodologiasAprendizagensNo2022}. Dados do Instituto Nacional de Estudos e Pesquisas Educacionais Anísio Teixeira (INEP) reforçam essa tendência ao apontar um crescimento de 378,9\% nas matrículas em cursos de graduação a distância no período de 2009 a 2019. \cite{moranEnsinoSuperiorDistancia2009} destaca que a relevância do EAD está associada não apenas à sua expansão numérica, mas também à capacidade de flexibilizar o processo de aprendizagem e aproximá-lo das necessidades da sociedade contemporânea.

Entretanto, mesmo com os avanços proporcionados pelas Tecnologias da Informação e Comunicação (TIC), há limitações importantes. Suhre Ribeiro (2010, \textit{apud} \cite{schneiderSalaAulaInvertida2013}) observa que, em muitos casos, ocorre apenas a transposição da sala de aula tradicional para o ambiente remoto, sem adaptação metodológica. Esse movimento pode comprometer o potencial do EAD, já que, como afirma Belloni (2002)\cite{belloniEnsaioSobreEducacao2002}, essa modalidade permite explorar a diversidade inerente aos grupos sociais e produzir experiências educativas mais efetivas. A aplicação de metodologias centradas na exposição, marcadas pela comunicação unilateral e pelo protagonismo do professor, pode, assim, representar uma limitação conceitual no contexto do ensino remoto. \cite{belloniEnsaioSobreEducacao2002} finaliza dizendo que "A EAD tem como principal característica reduzir distâncias entre docentes e discentes (bem como demais interagentes), sejam elas espaciais, culturais, sociais e temporais".

Entre os pontos fortes do EAD, destacam-se a democratização do acesso à educação, a flexibilidade, a autonomia e a possibilidade de inovação metodológica e tecnológica. Muller (2009, \textit{apud} \cite{vasconcelosEstrategiasEnsinoAplicaveis2013}) ressalta que "tecnologias de informação e comunicação, em convergência com as telecomunicações e redes, mais especificamente a \textit{internet}, podem se tornar um meio de excelência para promover educação continuada e a distância". No entanto, os desafios não podem ser ignorados: taxas elevadas de evasão, dispersão da atenção, dificuldades de acompanhamento individualizado e desigualdade no acesso às tecnologias configuram obstáculos recorrentes no contexto brasileiro.

Portanto, é essencial avaliar quais técnicas de ensino conseguem potencializar os recursos tecnológicos disponibilizados pela Quarta Revolução Industrial, de modo a garantir não apenas a ampliação do acesso, mas também a qualidade e a permanência dos estudantes na modalidade a distância.